%------------------------------------------------------------------------------
% Chapter 1: Introduction
%------------------------------------------------------------------------------
% Version: 2.0 (January 2026)
%
% Purpose: Introduces the Evidence-Based Framework for Economic and Social
%          Behavior (EBF). Establishes why integration is now possible and
%          presents the four core concepts (C, Ψ, C*, K) plus the 10C CORE preview.
%
% Primary Appendix: None (Foundation Chapter)
% Secondary Appendices: LIT-META (Metascience), G (Glossary), A (Limiting Cases),
%                       AU/AV/AW (Behavioral Pipeline)
%
% Prerequisites: None (entry point)
% Leads to: Chapter 1b (10C CORE Architecture), Chapter 2 (Rationality)
%
% Chapter Type: B (Foundation)
%------------------------------------------------------------------------------


%%%%%%%%%%%%%%%%%%%%%%%%%%%%%%%%%%%%%%%%%%%%%%%%%%%%%%%%%%%%%%%%%%%%%%%%%%%%%%%
% CHAPTER SCOPE BOX
%%%%%%%%%%%%%%%%%%%%%%%%%%%%%%%%%%%%%%%%%%%%%%%%%%%%%%%%%%%%%%%%%%%%%%%%%%%%%%%
\begin{tcolorbox}[
    colback=purple!5!white,
    colframe=purple!75!black,
    title={\textbf{Chapter Scope: Ziel / In-Scope / Out-of-Scope}},
    fonttitle=\bfseries
]

\textbf{Ziel (Objective):}
Develop the conceptual understanding of framework overview, motivation as a foundation for the EBF framework.

\vspace{0.3cm}
\textbf{In-Scope:}
\begin{itemize}[nosep]
    \item Conceptual introduction to framework overview
    \item Motivation and intuition
    \item Connection to subsequent chapters
    \item Key definitions and concepts
\end{itemize}

\vspace{0.3cm}
\textbf{Out-of-Scope (delegiert an Appendices):}
\begin{itemize}[nosep]
    \item Complete formal treatment $\rightarrow$ FORMAL appendices
    \item Empirical validation $\rightarrow$ METHOD appendices
    \item Domain applications $\rightarrow$ Application chapters
\end{itemize}

\end{tcolorbox}



%%%%%%%%%%%%%%%%%%%%%%%%%%%%%%%%%%%%%%%%%%%%%%%%%%%%%%%%%%%%%%%%%%%%%%%%%%%%%%%
% INTUITION BOX
%%%%%%%%%%%%%%%%%%%%%%%%%%%%%%%%%%%%%%%%%%%%%%%%%%%%%%%%%%%%%%%%%%%%%%%%%%%%%%%
\begin{tcolorbox}[
    colback=blue!5!white,
    colframe=blue!75!black,
    title={\textbf{Intuition: A Motivating Example}},
    fonttitle=\bfseries
]
Consider \textbf{Anna}, a professional facing a decision that involves framework overview, motivation.

She initially evaluates the choice using standard criteria, but something feels incomplete.
\textbf{Thomas}, her colleague, suggests considering additional dimensions beyond the obvious.

This chapter explores what Anna discovers when she applies the EBF framework---how
framework overview, motivation interact with broader welfare considerations.

\medskip
\textit{By the end of this chapter, you will understand why Anna's initial analysis
was incomplete and how the EBF approach provides a more comprehensive view.}
\end{tcolorbox}



\section{Introduction}
\label{sec:introduction}
%%%%%%%%%%%%%%%%%%%%%%%%%%%%%%%%%%%%%%%%%%%%%%%%%%%%%%%%%%%%%%%%%%%%%%%

%------------------------------------------------------------------------------
% QUICK REFERENCE BOX
%------------------------------------------------------------------------------
\begin{tcolorbox}[colback=gray!5!white, colframe=gray!75!black,
    title=\textbf{Begriffe in diesem Kapitel}]

\small
\begin{tabular}{@{}p{2.8cm}p{10cm}@{}}
\textbf{Complementarity} & Wie stark beeinflusst Verhalten einer Person den Nutzen anderer? $C_{ij} = \partial^2 U/\partial x_i \partial x_j$ \\[3pt]
\textbf{Context ($\Psi$)} & Die 8 Dimensionen, die ``Spielregeln'' bestimmen: Institutionen, Soziales, Kognitiv, Information, Ökonomisch, Temporal, Markt, Flexibilität \\[3pt]
\textbf{Reference ($C^*$)} & Das kontextabhängige ``ideale'' Verhaltensmuster; $C^*(\Psi)$ \\[3pt]
\textbf{Coherence ($K$)} & Wie gut passt Verhalten zur Referenz? $K = 1 - \|C-C^*\|/\|C^*\|$ \\[3pt]
\textbf{Metatheory} & Framework das organisiert, \emph{wann} welche Theorie gilt---ersetzt keine \\
\end{tabular}

\vspace{3pt}
\textit{Vollständiges Glossar mit intuitiven und formalen Definitionen: Appendix~G}
\end{tcolorbox}

%------------------------------------------------------------------------------
% APPENDIX REFERENCES BOX
%------------------------------------------------------------------------------
\begin{tcolorbox}[
    colback=yellow!10!white,
    colframe=orange!75!black,
    title={\textbf{Appendix References for This Chapter}},
    fonttitle=\bfseries
]
\small
\begin{tabular}{@{}p{3.5cm}p{5.5cm}p{3.5cm}@{}}
\textbf{Appendix} & \textbf{Content} & \textbf{Section} \\
\midrule
LIT-META (Metascience) & Why integration emerges now & \ref{sec:intro-history} \\
G (REF-GLOSSARY) & Symbol definitions & All sections \\
A (FORMAL-DERIVE) & Limiting cases (Arrow-Debreu) & \ref{sec:intro-concepts} \\
AU, AV, AW (CORE) & 10C Behavioral Pipeline & \ref{sec:intro-9c} \\
F (REF-EXAMPLES) & Worked examples & \ref{sec:intro-validation}
\end{tabular}
\end{tcolorbox}

%------------------------------------------------------------------------------
% INTUITION BOX
%------------------------------------------------------------------------------
\begin{tcolorbox}[colback=gray!5!white, colframe=gray!75!black,
    title=\textbf{Why This Matters: Anna's Dilemma}]
\small
Anna is a behavioral economist studying charitable giving. She runs an experiment in Switzerland: a matching grant doubles donations. Result: donations increase 45\%. She replicates in rural Peru: donations increase 8\%. Same intervention, different context.

\textbf{Classical economics} says: ``Estimate the average treatment effect.'' But which average---Swiss or Peruvian? Both?

\textbf{This framework} says: The matching grant works \emph{differently} depending on trust ($\Psi_S$), institutional enforcement ($\Psi_I$), and economic resources ($\Psi_E$). Switzerland has high values on all three; Peru's context differs. The effect isn't ``45\%'' or ``8\%''---it's $\tau(\Psi)$, a function of context.

\textit{This chapter explains how context ($\Psi$) transforms every economic relationship from a single number into a context-dependent function.}
\end{tcolorbox}

%------------------------------------------------------------------------------
% CENTRAL QUESTION BOX
%------------------------------------------------------------------------------
\paragraph{The Central Question.}
This chapter addresses the foundational question of the entire framework:

\begin{center}
\fcolorbox{blue!75!black}{blue!5!white}{\parbox{0.85\textwidth}{
\textbf{Central Question:} Why does economics appear fragmented, and what would it take to create an integrative framework that organizes rather than replaces existing theories?
}}
\end{center}

%------------------------------------------------------------------------------
% CHAPTER OVERVIEW
%------------------------------------------------------------------------------
\paragraph{Chapter Overview.}
\label{sec:intro-overview}
This chapter proceeds as follows:
\begin{itemize}
    \item Section \ref{sec:intro-problem}: The problem of disciplinary fragmentation
    \item Section \ref{sec:intro-history}: The historical path dependency that created fragmentation
    \item Section \ref{sec:intro-question}: When integration becomes possible
    \item Section \ref{sec:intro-metatheory}: Our solution---metatheory through complementarity
    \item Section \ref{sec:intro-concepts}: The four core concepts ($C$, $\Psi$, $C^*$, $K$)
    \item Section \ref{sec:intro-9c}: Preview of the 10C CORE Framework
    \item Section \ref{sec:intro-validation}: Empirical validation
    \item Section \ref{sec:intro-methodology}: Methodological distinction from LLM approaches
    \item Section \ref{sec:intro-roadmap}: Roadmap for the remainder of the paper
\end{itemize}

\subsection{The Problem: Fragmentation}
\label{sec:intro-problem}

Economics appears fragmented. Classical theory explains market efficiency but struggles with social preferences. Behavioral economics explains social preferences but often holds institutions fixed. Institutional economics models institutional change but frequently ignores micro-level dynamics. Finance explains financial crises but rarely traces political consequences.

Each subdiscipline has its own assumptions, methods, and vocabulary. They coexist uneasily, sometimes contradicting each other, often talking past each other. Graduate students learn them as separate courses. Policy makers receive conflicting advice.

This fragmentation is not merely inconvenient---it is intellectually unsatisfying and practically costly. The great challenges of our time---climate change, inequality \citep{piketty2014}, institutional decay, technological disruption---do not respect disciplinary boundaries. They involve interactions between markets, behaviors, institutions, and identities. Understanding them requires a framework that integrates rather than fragments.\footnote{Appendix~LIT-META provides a metascientific analysis of why this fragmentation persists and why integration has been systematically underproduced by academic institutions.}

\textit{Why has integration not occurred despite universal recognition of its value?
Appendix~LIT-META provides the metascientific answer: five structural barriers---academic
incentives, disciplinary silos, peer review bias, funding structures, and territorial
defense of paradigms---systematically discourage integration while rewarding fragmentation.
The barrier is institutional, not intellectual.}

\subsection{The History: Economics and Technological Possibility}
\label{sec:intro-history}

To understand why economics appears fragmented, we must understand how it developed---and why that development was shaped by technological constraints that no longer apply.

\paragraph{The Path Dependency.}
The history of modern economics follows a clear sequence:
\begin{enumerate}
    \item \textbf{1947:} Paul Samuelson's \emph{Foundations of Economic Analysis} created a mathematical \emph{framework}---a language for expressing economic relationships \citep{samuelson1947foundations}.
    \item \textbf{1953:} Milton Friedman's ``as if'' methodology defended \emph{unrealistic assumptions} if they yielded good predictions \citep{friedman1953essays}.
    \item \textbf{1959:} Arrow and Debreu axiomatized general equilibrium under separability ($C_{ij} = 0$)---the only case analytically solvable with pencil and paper \citep{arrowdeb}.
    \item \textbf{1970s:} The Chicago School institutionalized this methodology \citep{stigler1982,becker1976economic}, treating the framework as an empirically validated model.
    \item \textbf{1980s--2000s:} Behavioral anomalies accumulated (Kahneman-Tversky 1979, Güth 1982, Fehr-Schmidt 1999), but no integration was possible.
    \item \textbf{2010s:} The credibility revolution \citep{angrist2010,leamer2010tantalus} brought data-driven methods, yet theoretical fragmentation persisted.
    \item \textbf{2020s:} AI and modern computation finally enable what was always conceptually possible.
\end{enumerate}

\paragraph{Why Separability Was Necessary.}
The separability assumption---that my utility depends on my consumption, not yours; that your choices don't affect my optimal choice except through prices---was not merely convenient. It was \emph{computationally necessary}. Without computers, models with interacting agents ($C_{ij} \neq 0$) were analytically unsolvable. The Chicago methodology was a rational response to technological constraints, not an epistemological error.

Under separability, each agent's optimization problem is independent. Markets aggregate independent choices. Equilibrium is efficient. This Arrow-Debreu world was mathematically elegant, normatively powerful---and the only option given the tools available.

\paragraph{The Anomalies That Accumulated.}
Empirical research documented systematic violations:
\begin{itemize}
  \item \textbf{Social preferences:} People reject unfair offers even at personal cost \citep{guth1982,fehrschmidt}
  \item \textbf{Identity:} People sacrifice material welfare to maintain self-image \citep{akerlofkranton,benaboutirole}
  \item \textbf{Reference dependence:} Choices depend on context, framing, and expectations \citep{kahnemantversky1979}
  \item \textbf{Institutional effects:} The same people behave differently under different rules \citep{north1990institutions, acemoglu}
\end{itemize}

Each finding generated its own literature---behavioral economics \citep{camloewrabin}, identity economics, institutional economics. These literatures developed independently because \emph{no integration was technologically feasible}---and because scientific incentives favor specialization over synthesis \citep{smaldino2016natural,kuhn1962structure}. Without computation, context-dependent models with $C_{ij} \neq 0$ could not be operationalized.

\paragraph{Learning from All Social Sciences.}
Each social science developed methods suited to its questions and constraints. Psychology developed rich experimental paradigms for individual cognition. Sociology mapped social structures and cultural dynamics. Political science analyzed institutional design and collective action. Management science accumulated practical knowledge about organizational change. Each contributed insights that this framework integrates.

Economics contributed something distinctive: a common mathematical language that enabled cumulative development \citep{fourcade2009economists,fourcade2015}. This was not superiority but \emph{adaptation}---economics developed the formal tools that were computationally tractable at each epoch. The result was a discipline that could absorb new methods (experiments, field studies, machine learning) while maintaining theoretical coherence.

\paragraph{Why Integration Must Come from Economics.}
This framework can only emerge from the economic tradition---not because economists are smarter, but because economics built the necessary infrastructure. Seventy years of formalization created something no other social science possesses: a mathematical backbone onto which new insights can be integrated.

Consider what integration requires: (1) a common formal language to express relationships precisely, (2) cumulative theoretical development to build upon, and (3) methodological flexibility to absorb new evidence from any source. Now consider each social science:
\begin{itemize}
    \item \textbf{Psychology} developed rich experimental insights into individual cognition \citep{kahneman2011,stanovich2011rationality}, but Freudians, behaviorists, cognitive psychologists, and neuroscientists use incompatible frameworks. No common language exists for integration \citep{mischel2008toothbrush}.
    \item \textbf{Sociology} mapped social structures and cultural dynamics \citep{merton1968social,giddens1984constitution}, but explicitly rejected formalization. Functionalism, conflict theory, and symbolic interactionism coexist without any integration mechanism \citep{collins1994four}.
    \item \textbf{Political science} analyzed institutional design and collective action \citep{ostrom1990,olson1965logic}, but only partially formalized. Rational choice exists alongside historical institutionalism without synthesis \citep{hall1996political}.
    \item \textbf{Management science} accumulated practical knowledge about organizational change \citep{lewin1947frontiers,kotter1996leading}, but exhibits fashion cycles rather than cumulation \citep{abrahamson1996management}. Each decade's framework replaces rather than builds on the previous.
    \item \textbf{Economics} developed all three: a common mathematical language (utility, equilibrium, optimization) \citep{mascolell1995microeconomic}, cumulative development (each generation builds on predecessors), and methodological flexibility (absorbed experiments, RCTs, machine learning while maintaining coherence) \citep{angrist2009mostly}.
\end{itemize}

Psychology's experimental insights, sociology's structural analysis, heterodox economics' contextual sensitivity---all can be integrated, but only onto a formal structure that can express them precisely. That structure exists only in economics.

This is why the timing matters. The formal infrastructure had to exist \emph{before} AI could operationalize it. What Samuelson built, what Arrow-Debreu axiomatized, what behavioral economics tested---this seventy-year accumulation is the foundation. AI does not replace it; AI finally makes it usable for the non-separable, context-dependent models that were always conceptually possible but computationally infeasible. The synthesis is collaborative: other social sciences contribute content; economics contributes structure. Appendix~LIT-META documents this infrastructure thesis in detail.

\paragraph{Building on Heterodox Insights.}
Heterodox economics---Post-Keynesian \citep{lavoie2014post}, Institutionalist \citep{hodgson2004evolution}, Austrian \citep{hayek1945knowledge}, and other traditions---identified crucial phenomena that separable models could not capture: institutional path dependence, fundamental uncertainty, the social construction of preferences \citep{lawson2006nature,davis2008heterodox}. These insights were correct and important.

What was missing was not critique but \emph{tools}. As Colander observed, mainstream economics eventually formalized the very topics heterodox economists had pioneered \citep{colander2004changing}. The insights were absorbed; the question was always \emph{how} to operationalize them.

This framework builds on heterodox insights by providing the formal structure they lacked. Context ($\Psi$) captures the institutional and cultural factors heterodox economists emphasized. Non-separability ($C_{ij} \neq 0$) formalizes the social interdependencies they described. The framework does not replace heterodox contributions---it gives them mathematical expression that modern computation can operationalize. Appendix~LIT-META documents this synthesis in detail.

\subsection{The Question: When Does Integration Become Possible?}
\label{sec:intro-question}

The question is not ``why has integration failed?''---integration was not possible given technological constraints. The question is: \emph{when do conditions permit an integrative framework?}

\paragraph{Four Conditions for Integration.}
Integration emerges naturally when:
\begin{enumerate}
    \item Sufficient anomalies have accumulated to reveal the common pattern
    \item A framework (rather than another competing model) emerges that can organize them
    \item The framework approach avoids zero-sum competition with existing literatures
    \item Technology permits operationalization of the framework
\end{enumerate}

\paragraph{All Four Conditions Are Now Satisfied.}
After fifty years of behavioral economics, conditions (1)--(3) have long been met. The evidence base is rich. The pattern---context-dependent complementarity---is clear. The framework approach allows existing theories to keep their insights.

What changed is condition (4): modern computation and AI now permit operationalization. Models with $C_{ij} \neq 0$ can be solved numerically. Context ($\Psi$) can be estimated using LLM Monte Carlo methods (Appendix~AN). The framework-model distinction can finally be maintained operationally: the framework organizes thinking, while situation-specific models are generated and tested.

\paragraph{Looking Forward: The Next Fifteen Years.}
This technological moment is transformative. The same 8$\Psi$ framework that organizes economic behavior can be applied to \emph{the development of economics itself}. Appendix~LIT-META develops this model in detail, projecting an inflection point around 2035 when AI-powered frameworks demonstrably outperform traditional separable models in prediction accuracy.

Quality control shifts from pre-publication peer review to post-publication replication and prediction success. The theory-practice separation that defined twentieth-century academia dissolves when empirical performance becomes the arbiter of scientific validity. This is not speculation---it is a model with testable predictions that the next fifteen years will confirm or refute.

\subsection{Our Solution: Metatheory Through Complementarity}
\label{sec:intro-metatheory}

This paper provides the missing framework. But we make a crucial distinction: our contribution is not another \emph{theory} competing with existing ones, but a proposed \emph{metatheory} that aims to organize them.

\textit{This distinction is historically decisive. As Appendix~LIT-META documents,
Samuelson's \emph{Foundations} (1947) was conceived as a framework---a mathematical
language for economics---but was treated as a model with testable predictions. When
anomalies accumulated, the ``model'' was falsified but no framework remained to integrate
them. Result: fragmentation. EBF is explicitly a framework that generates testable,
situation-specific models while remaining unfalsifiable at the meta-level.}

\paragraph{What a Metatheory Does.}
A metatheory does not replace component theories---it specifies when each applies. Consider an analogy: Newtonian mechanics, relativistic mechanics, and quantum mechanics are distinct theories. But physics also has organizing principles (symmetry, conservation, least action) that explain \emph{when} each theory applies and \emph{how} they relate.

Our framework provides this organizing structure for economics:

\begin{itemize}
    \item \textbf{Arrow-Debreu} applies when complementarity is negligible ($C_{ij} \approx 0$) and context is stable ($\dot{\Psi} \approx 0$)
    \item \textbf{Behavioral economics} applies when individual biases dominate ($\Psi_C$ varies) but institutions are fixed
    \item \textbf{Institutional economics} applies when rules dominate ($\Psi_I$ varies) but individual behavior aggregates
    \item \textbf{Our full framework} applies when multiple $\Psi$ dimensions vary simultaneously with significant complementarity
\end{itemize}

The decisive advantage: existing theories tell you \emph{that} certain effects occur. Our framework tells you \emph{when} they occur---which contexts amplify, dampen, or reverse each effect.

\subsection{The Four Core Concepts}
\label{sec:intro-concepts}

This section introduces the four foundational concepts of the framework. Each concept is defined precisely here; Appendix~G provides a comprehensive glossary with both textual explanations and mathematical formalizations for all technical terms used throughout this paper.

\paragraph{1. Complementarity Matrix $C$.}
The second derivatives of utility functions:
\begin{equation}
  C_{ij} = \frac{\partial^2 U_i}{\partial x_i \partial x_j}
\end{equation}

This captures how my marginal utility changes when you change your behavior.
\begin{itemize}
  \item $C_{ij} > 0$: Complements---your cooperation makes my cooperation more valuable
  \item $C_{ij} < 0$: Substitutes---your action crowds out mine
  \item $C_{ij} = 0$: Independent---the Arrow-Debreu case (see Appendix~A for formal limiting cases)
\end{itemize}

\paragraph{2. Context Vector $\Psi \in \mathbb{R}^8$.}
Technology, institutions, norms, and culture---everything that shapes the payoff structure but is not under any single agent's control.

We operationalize $\Psi$ through eight primitive dimensions (Table~\ref{tab:psi-intro}), each derived from decision-theoretic requirements and independently measurable with established instruments.

\begin{table}[htbp]
\centering
\caption{The Eight Context Dimensions}
\label{tab:psi-intro}
\small
\begin{tabular}{@{}clll@{}}
\toprule
\textbf{Symbol} & \textbf{Dimension} & \textbf{What It Captures} & \textbf{Key Measures} \\
\midrule
$\Psi_I$ & Institutional & Formal rules and enforcement & Rule of Law, Polity IV \\
$\Psi_S$ & Social & Trust and social capital & WVS Trust, Civic Index \\
$\Psi_C$ & Cognitive & Information processing capacity & PISA, Financial Literacy \\
$\Psi_K$ & Informational & Transparency and signal quality & Press Freedom, CPI \\
$\Psi_E$ & Economic & Development and resources & GDP/capita, Credit Access \\
$\Psi_T$ & Temporal & Stability and time horizons & EPU Index, Inflation Vol. \\
$\Psi_M$ & Market Scope & Geographic reach and openness & Trade/GDP, LPI \\
$\Psi_F$ & Factor Flexibility & Adjustment capacity & EPL Index, Mobility \\
\bottomrule
\end{tabular}
\end{table}

\paragraph{Empirical Grounding of Each Dimension.}
Each $\Psi$ dimension is grounded in established research programs:
\begin{itemize}
  \item \textbf{$\Psi_C$ (Cognitive):} Bounded rationality \citep{simon1962},
        prospect theory \citep{kahnemantversky1979}, dual-process theory
        \citep{kahneman2011}. See Appendix~U for Kahneman-Tversky foundations.
  \item \textbf{$\Psi_I$ (Institutional):} Institutional economics \citep{north1990institutions},
        political institutions \citep{acemoglurobinson2012}. See Appendix~L.
  \item \textbf{$\Psi_S$ (Social):} Social preferences including inequality aversion,
        reciprocity, and image concerns \citep{fehrcharness2024}. See Appendix~AJ.
  \item \textbf{$\Psi_K$ (Informational):} Information economics
        \citep{akerlof1970,spence1973,stiglitz2000}. See Appendix~W.
  \item \textbf{$\Psi_T$ (Temporal):} Time preferences and present bias
        \citep{laibson1997,frederick2002time}. See Appendix~AH.
  \item \textbf{$\Psi_E$ (Economic):} Resource constraints, budget effects,
        capital markets \citep{fama1970efficient}. See Appendix~Y.
  \item \textbf{$\Psi_M$ / $\Psi_P$ (Market/Physical):} Economic geography \citep{krugman1991},
        urban economics \citep{glaeser2008}. See Appendix~AI.
  \item \textbf{$\Psi_F$ / $\Psi_{Cult}$ (Flexibility/Cultural):} Cultural dimensions 
        \citep{hofstede1980,house2004}, world values \citep{inglehart2014}. See Appendix~V.
\end{itemize}

The Global Preference Survey \citep{falk2018} provides cross-country measures
for patience ($\Psi_T$), risk tolerance, social preferences ($\Psi_S$), and
trust across 76 countries, demonstrating substantial and systematic variation
that predicts economic outcomes.


Crucially, $\Psi$ is \emph{endogenous}:
\begin{equation}
  \Psi_{t+1} = f(\Psi_t, a_t)
\end{equation}
Actions today shape the context tomorrow. This feedback loop creates path dependence, institutional persistence, and the possibility of regime change.

\paragraph{3. Reference Structure $C^*(\Psi)$.}
Given context $\Psi$, there exists a \emph{self-consistent} complementarity structure---the configuration where beliefs about others' behavior match actual behavior. This is the equilibrium toward which systems gravitate.

\paragraph{4. Coherence Index $K$.}
The alignment between actual structure $C$ and reference structure $C^*$:
\begin{equation}
    K = 1 - \frac{\|C - C^*(\Psi)\|}{\|C^*(\Psi)\|}
\end{equation}
High coherence ($K \approx 1$) means stability. Low coherence ($K \ll 1$) means pressure for change. Crucially, coherence is distinct from quality: a system can be coherent but bad (stable dystopia) or incoherent but improving (beneficial transition).

\paragraph{From Four Concepts to Eight Questions.}
\label{sec:intro-9c}
These four core concepts---Complementarity ($C$), Context ($\Psi$), Reference Structure ($C^*$), and Coherence ($K$)---answer \emph{how interactions work} and \emph{when context matters}. But they leave critical questions unanswered: Who exactly has utility? What does utility consist of? Where do parameters come from? And crucially: what transforms utility into action?

The \textbf{10C CORE Framework} completes this architecture by addressing these gaps. Each CORE question has a dedicated appendix with 80--150 axioms, complete formalization, and empirical grounding.

%==============================================================================
% 10C CORE PREVIEW BOX
%==============================================================================
\begin{tcolorbox}[colback=blue!5!white, colframe=blue!75!black,
    title=\textbf{The 10C CORE Framework}]

These four concepts are the foundation. But complete behavioral prediction
requires answering \textbf{eight fundamental questions}:

\begin{center}
\small
\begin{tabular}{@{}cllll@{}}
\toprule
\textbf{\#} & \textbf{CORE} & \textbf{Question} & \textbf{Output} & \textbf{Appendix} \\
\midrule
1 & WHO & Who has utility? & Level $L \in \{0...\infty\}$ & AAA \\
2 & WHAT & What is utility? & Dimensions $d$ (FEPSDE) & C \\
3 & HOW & How do they interact? & Complementarity $\gamma$ & B \\
4 & WHEN & When does context matter? & Context $\Psi$ (8 dim) & V \\
5 & WHERE & Where do numbers come from? & Parameters $\Theta$ & BBB \\
6 & AWARE & How salient is utility? & Awareness $A(\cdot)$ & AU \\
7 & READY & How willing to act? & Willingness $WAX, \theta$ & AV \\
8 & STAGE & Where in change journey? & Stage $S(t)$ & AW \\
\bottomrule
\end{tabular}
\end{center}

\textbf{The 4-Stage Pipeline:}
\begin{enumerate}[nosep]
    \item \textbf{Utility Definition (5C):} $U_{pot} = f(L, d, \gamma, \Psi, \Theta)$
    \item \textbf{Awareness Filter (6th C):} $U_{eff} = A(\cdot) \times U_{pot}$
    \item \textbf{Action Threshold (7th C):} Action $\Leftrightarrow WAX \geq \theta$
    \item \textbf{Behavioral Change (8th C):} $S(t+1) = f(S(t), A, WAX, \theta, \Psi)$
\end{enumerate}

\textbf{Key Insight:} The four concepts above (Complementarity, Context, Reference, Coherence) answer questions 3--4 (HOW and WHEN). The 10C Framework extends this to the complete behavioral prediction pipeline---from potential utility through sustained behavioral change.

\textit{See Chapter~1b for the complete 10C CORE architecture.}

\end{tcolorbox}

\subsection{Empirical Validation: The Eight $\Psi$ Dimensions}
\label{sec:intro-validation}

The framework is not merely theoretical. We test whether the eight $\Psi$ dimensions actually explain treatment effect heterogeneity across contexts.

\paragraph{The Core Prediction.}
If context matters as we claim, then:
\begin{equation}
    \tau(X \rightarrow Y | \Psi) = \beta_0 + \sum_{k=1}^{8} \beta_k \Psi_k + \text{interactions}
\end{equation}
Treatment effects should be systematically related to context dimensions.

\paragraph{Results.}
Across six canonical economic relationships, the eight dimensions explain substantial heterogeneity:

\begin{table}[htbp]
\centering
\caption{Empirical Validation: 8$\Psi$ Model Across Six Outcomes}
\label{tab:validation-intro}
\small
\begin{tabular}{@{}lccc@{}}
\toprule
\textbf{Relationship} & \textbf{$R^2$} & \textbf{Adj. $R^2$} & \textbf{Dominant $\Psi$} \\
\midrule
Education $\rightarrow$ Earnings & 84.5\% & 76.8\% & $\Psi_K$, $\Psi_M$ \\
Management $\rightarrow$ Productivity & 78.9\% & 68.3\% & $\Psi_F$ \\
Democracy $\rightarrow$ Growth & 39.8\% & 9.7\% & $\Psi_I$, $\Psi_E$ \\
Patience $\rightarrow$ Growth & 67.3\% & 51.0\% & $\Psi_T$, $\Psi_K$ \\
Information $\rightarrow$ Behavior & 82.4\% & 73.7\% & $\Psi_K$ (--), $\Psi_F$ \\
Income $\rightarrow$ Health & 67.7\% & 51.6\% & $\Psi_T$, $\Psi_I$ \\
\midrule
\textbf{Average} & \textbf{70.1\%} & \textbf{55.2\%} & --- \\
\bottomrule
\end{tabular}
\end{table}

Three findings emerge:
\begin{enumerate}
    \item \textbf{Different dimensions dominate for different effects.} Factor flexibility ($\Psi_F$) matters most for management; information quality ($\Psi_K$) for education returns; temporal stability ($\Psi_T$) for patience effects.
    
    \item \textbf{Micro effects are better explained than macro effects.} Individual decisions show more systematic context-dependence ($R^2 \approx 75$--$85\%$) than aggregate outcomes ($R^2 \approx 40$--$70\%$).
    
    \item \textbf{No ninth dimension emerged.} We tested ethnic fractionalization, natural resources, and implementation capacity. None improved the model across all outcomes (see Appendix~\ref{app:psi} for methodology).
\end{enumerate}

\subsection{The Framework's Falsifiable Claim}
\label{sec:intro-falsifiable}

What would refute this framework? We commit to a specific, testable proposition:

\begin{quote}
\textbf{Core Falsifiable Claim:} \emph{No economic effect is context-free.}
\end{quote}

If any intervention produces identical effects across all $\Psi$ configurations---if education returns are the same in Switzerland and Nigeria, if information campaigns work equally in transparent and opaque environments, if management practices improve productivity identically with rigid and flexible labor markets---then our framework would face serious empirical challenges.

This is a strong hypothesis. Standard theory often assumes context-independence (the ``average treatment effect'' literature). We hypothesize that context-dependence is universal and systematic.

\subsection{Methodological Distinction: A Self-Improving Research Architecture}
\label{sec:intro-methodology}

A natural question arises: why develop a new framework when recent advances in artificial intelligence offer an alternative approach? Large language models (LLMs) can be conditioned on demographic information to create ``digital twins'' that simulate human responses \citep{horton2023,argyle2023}. Could these models replace or complement human experiments?

The answer appears to be no. \citet{peng2025} conducted 19 preregistered studies comparing a representative U.S. panel with their digital twins---LLMs conditioned on over 500 variables per individual. The results were sobering: digital twins replicated only 50\% of classic behavioral economics experiments, often reproducing \emph{textbook predictions that humans violate}.

This failure illuminates a fundamental distinction:
\begin{itemize}
    \item \textbf{LLMs} train on text \emph{about} experiments. They learn that ``the endowment effect is a robust phenomenon.'' They do not learn \emph{when} it is strong versus weak, \emph{how} it varies with context, or \emph{why} it sometimes disappears entirely.
    
    \item \textbf{This framework} calibrates on experimental \emph{data themselves}. Fifty years of behavioral economics provides not illustration but calibration data for $C^*(\Psi)$---the context-dependent reference structure.
\end{itemize}

The distinction is between \emph{propositional knowledge} (``effect X exists'') and \emph{functional knowledge} ($Y = f(X, \Psi)$ with estimated parameters). LLMs excel at the former; prediction requires the latter.

Table~\ref{tab:llm-comparison} summarizes the key differences.

\begin{table}[htbp]
\centering
\caption{Comparison: LLM-Based Simulation vs.\ Behavioral Calibration}
\label{tab:llm-comparison}
\small
\begin{tabular}{@{}lcc@{}}
\toprule
\textbf{Aspect} & \textbf{LLM / Digital Twin} & \textbf{This Framework} \\
\midrule
Data source & Text corpora & Experimental behavioral data \\
Knowledge type & Propositional (``X exists'') & Functional ($Y = f(X, \Psi)$) \\
Heterogeneity & Averaged away & Explicitly modeled via $\Psi$ \\
Replication rate & $\sim$50\% & $>$85\% (by construction) \\
Cumulation & Requires retraining & Each experiment improves $C^*$ \\
Self-improvement & None (static) & Deviations trigger revision \\
Time to incorporate new evidence & Months (retraining) & Days (parameter update) \\
\bottomrule
\end{tabular}
\end{table}

\paragraph{Worked Example: Cross-Cultural Ultimatum Game.}
Consider predicting rejection rates in the ultimatum game across contexts. An LLM, asked to predict behavior in the Machiguenga (Peru) versus Pittsburgh, might respond: ``Rejection rates vary across cultures, with Western populations typically showing higher rejection of unfair offers due to stronger fairness norms.'' This is propositional knowledge---qualitatively correct but quantitatively useless.

The present framework, calibrated on \citet{henrich2001} and subsequent cross-cultural studies, generates:
\begin{equation}
    \text{Rejection}_{20\%}(\Psi) = 0.12 + 0.35 \cdot \Psi_S + 0.18 \cdot \Psi_I - 0.22 \cdot \Psi_E
\end{equation}
For Machiguenga ($\Psi_S = 0.3$, $\Psi_I = 0.2$, $\Psi_E = 0.1$): predicted rejection $\approx 0.25$ (actual: 0.05--0.25). For Pittsburgh ($\Psi_S = 0.7$, $\Psi_I = 0.8$, $\Psi_E = 0.9$): predicted rejection $\approx 0.38$ (actual: 0.40--0.50). The framework captures not just \emph{that} behavior varies but \emph{how much} and \emph{why}---because it was calibrated on the variation itself.

\paragraph{Cumulative Learning.}
Beyond accuracy, a crucial difference concerns \emph{cumulation}. When a new behavioral experiment is published:
\begin{itemize}
    \item \textbf{For LLMs:} The new evidence cannot be incorporated without retraining. The model remains frozen at its last training cutoff.
    \item \textbf{For this framework:} The new experiment becomes another data point for $C^*(\Psi)$. Each experiment improves the reference structure---precision increases with the calibration base.
\end{itemize}

Section~4.X develops this argument in detail, showing how predictive accuracy scales with the number of calibrating experiments: 50 experiments yield main effects ($R^2 \approx 0.3$--$0.5$), 500 experiments reveal interactions ($R^2 \approx 0.5$--$0.7$), and 1,500+ experiments capture the full complementarity structure ($R^2 \approx 0.7$--$0.85$). This is not aspirational but a logical consequence of the information content in behavioral variation.

\paragraph{Self-Reinforcement Learning.}
The framework goes beyond passive cumulation. When predictions deviate from observations ($\Delta = C - C^*(\Psi) \neq 0$), these deviations are not noise to be averaged away but \emph{informative signals} that trigger structured revision. Appendix~AK develops a five-type diagnostic protocol that classifies deviations (measurement error, missing moderator, context misclassification, structural inadequacy, genuine anomaly), and Appendix~AL shows how each type translates into specific model updates.

This creates a feedback loop analogous to:
\begin{itemize}
    \item \textbf{AlphaZero}: Self-play generates data that improves the policy, which generates better data
    \item \textbf{Numerical weather prediction}: Forecast errors update model parameters, improving tomorrow's forecast
\end{itemize}

The result is not a static model but a \emph{self-improving research architecture}---one that accelerates the traditional theory-test-revise cycle from decades to months.

The irony is striking: LLMs often reproduce the rational-choice predictions that behavioral economics was invented to correct. A behavioral prediction system should capture \emph{when and why} people deviate from rationality---not reproduce the deviations' textbook descriptions.

\subsection{Roadmap}
\label{sec:intro-roadmap}

The paper proceeds as follows:

\begin{itemize}
    \item \textbf{Chapter 1b} introduces the complete \textbf{10C CORE Framework}---the eight fundamental questions that structure behavioral prediction.

    \item \textbf{Section 2} reviews rationality concepts in classical economics.

    \item \textbf{Section 3} examines limits of utility maximization under stable preferences.

    \item \textbf{Section 4} establishes empirical foundations from theory to measurement, including Section~4.X which develops the methodological distinction between behavioral calibration and LLM-based simulation.

    \item \textbf{Section 5} develops complementarity as a structural principle (\textbf{CORE-HOW}).

    \item \textbf{Section 6} derives the reference structure $C^*(\Psi)$ through three approaches.

    \item \textbf{Section 7} analyzes complementarity, fit, and non-concavity.

    \item \textbf{Section 8} provides mathematical formalization.

    \item \textbf{Section 9} develops context as an endogenous variable (\textbf{CORE-WHEN}).

    \item \textbf{Section 10} presents the FEPSDE welfare framework (\textbf{CORE-WHAT}, \textbf{CORE-WHO}).

    \item \textbf{Section 11} develops the Awareness function---the salience filter that transforms potential into effective utility (\textbf{CORE-AWARE}).

    \item \textbf{Section 12} develops the Willingness function---the threshold mechanism that determines action (\textbf{CORE-READY}).

    \item \textbf{Section 13} introduces the Behavioral Change Journey---the temporal dimension of behavioral change (\textbf{CORE-STAGE}).

    \item \textbf{Section 14} synthesizes Willingness to Effectively Change along the journey.

    \item \textbf{Section 15} concludes.
\end{itemize}

\paragraph{Appendices.}
The 56 appendices are organized into 8 categories:

\begin{itemize}
    \item \textbf{CORE (10C Framework):} AAA (WHO---Level Hierarchy), B (HOW---Complementarity), C (WHAT---FEPSDE Matrix), V (WHEN---Context $\Psi$), BBB (WHERE---Parameter Estimation), AU (AWARE---Awareness Function), AV (READY---Willingness Function), AW (STAGE---Behavioral Change Journey)

    \item \textbf{FORMAL:} A (Theoretical Limiting Cases), D (Mathematical Proofs)

    \item \textbf{DOMAIN:} AA--AG (Labor, Matching, IO, Games, Mechanism Design, Social Choice, Complexity), W--Z (Information, Milgrom-Roberts, Capital Markets, Growth), AJ--AK (Social Preferences, Epistemics)

    \item \textbf{CONTEXT:} AH (Temporal), AI (Spatial)

    \item \textbf{METHOD:} AN (LLM Monte Carlo), AL (Self-Reinforcement), E (Empirical Operationalization), R (Validation)

    \item \textbf{PREDICT:} S (Falsifiable Predictions), AO--AT (Domain-Specific Predictions)

    \item \textbf{LIT:} I--Q (Nobel Contributions, Fehr, Acemoglu, Shleifer, Heckman, Autor, Duflo, Bloom), U (Kahneman-Tversky), LIT-META (Metascience Integration)

    \item \textbf{REF:} F (Worked Examples), G (Glossary), H (Computational History), T (Metatheory Response)
\end{itemize}

%------------------------------------------------------------------------------
% RELATED PUBLICATIONS BOX
%------------------------------------------------------------------------------
\begin{tcolorbox}[colback=purple!5!white, colframe=purple!75!black,
    title=\textbf{Related Publications}]
\small

This chapter has been adapted into a standalone working paper for academic circulation:

\begin{itemize}[nosep]
    \item \textbf{WP-2026-01:} Fehr, G. \& Beatrix Research Team (2026). ``Context-Dependent Complementarity: A Metatheoretical Framework for Integrating Behavioral Economics.'' FehrAdvice \& Partners AG Working Paper Series. \\
    \textit{Available at:} \texttt{papers/WP\_2026\_01\_Context\_Complementarity.pdf}
\end{itemize}

\textit{The working paper presents the core concepts of this chapter in a format suitable for peer review and academic citation.}
\end{tcolorbox}

%------------------------------------------------------------------------------
% SUMMARY
%------------------------------------------------------------------------------
\section{Summary}
\label{sec:ch1-summary}

\begin{tcolorbox}[
    colback=blue!5!white,
    colframe=blue!75!black,
    title={\textbf{Chapter 1 Summary: Introduction to the EBF Framework}}
]

\textbf{Core Insight:} The EBF framework integrates 50+ years of behavioral economics research under a unified structure based on complementarity ($\gamma$) and context ($\Psi$).

\textbf{Key Results:}
\begin{enumerate}[nosep]
    \item \textbf{10C CORE Architecture:} Eight fundamental questions (WHO, WHAT, HOW, WHEN, WHERE, AWARE, READY, STAGE) organize all behavioral phenomena
    \item \textbf{Complementarity Principle:} Utility dimensions interact ($\gamma \neq 0$), not just aggregate
    \item \textbf{Context Endogeneity:} Context $\Psi$ modulates all behavioral parameters
    \item \textbf{Behavioral Pipeline:} Awareness $\rightarrow$ Willingness $\rightarrow$ Action stages
\end{enumerate}

\textbf{Transition:} Chapter 2 examines how classical economics treated rationality and why extensions became necessary.
\end{tcolorbox}

\vspace{1em}

%------------------------------------------------------------------------------
% READING PATH BOX
%------------------------------------------------------------------------------
\begin{tcolorbox}[colback=green!5!white, colframe=green!75!black,
    title=\textbf{Reading Paths from Chapter 1}]
\small

\textbf{Continue reading:}
\begin{itemize}[nosep]
    \item \textbf{Chapter 1b:} Complete 10C CORE architecture---the eight fundamental questions
    \item \textbf{Chapter 2:} Rationality concepts in classical economics
\end{itemize}

\medskip
\textbf{Deep dive into specific topics:}
\begin{itemize}[nosep]
    \item \textbf{Why integration now?} Appendix UNMAPPED_MSC (Metascience Integration)
    \item \textbf{Arrow-Debreu as limiting case:} Appendix UNMAPPED_DER (FORMAL-DERIVE)
    \item \textbf{134 established theories catalog:} Appendix UNMAPPED_MS (REF-MODELSPACE)
    \item \textbf{Behavioral pipeline:} Appendices AU, AV, AW (Awareness, Willingness, Stage)
    \item \textbf{All terminology:} Appendix UNMAPPED_GLS (REF-GLOSSARY)
    \item \textbf{Worked examples:} Appendix UNMAPPED_EXM (REF-EXAMPLES)
\end{itemize}

\medskip
\textbf{Application-focused readers:}
\begin{itemize}[nosep]
    \item Skip to Chapter 14 (Behavioral Change Segments) for practical applications
    \item See Section~\ref{sec:intro-validation} for empirical validation summary
\end{itemize}
\end{tcolorbox}

%%%%%%%%%%%%%%%%%%%%%%%%%%%%%%%%%%%%%%%%%%%%%%%%%%%%%%%%%%%%%%%%%%%%%%%%%%%%%%%
% END OF CHAPTER 1: INTRODUCTION
%%%%%%%%%%%%%%%%%%%%%%%%%%%%%%%%%%%%%%%%%%%%%%%%%%%%%%%%%%%%%%%%%%%%%%%%%%%%%%%
